\documentclass[10pt,a4paper]{amsart}
\usepackage[margin=19mm]{geometry}
\usepackage{amsmath,amssymb,mathtools}
\usepackage[scr=rsfs,cal=euler]{mathalfa}
\usepackage{xfrac}
\usepackage[unicode,hidelinks,bookmarks=false]{hyperref}
\newcommand{\ie}{i.e.\ }
\newcommand{\cf}{cf.\ }
\newcommand{\resp}{respectively\ }
\newcommand{\Subsection}{\subsection}
\providecommand{\nobreakdash}{\nobreak\hbox{-}}
\setlength{\emergencystretch}{2em}
\multlinegap0pt
\usepackage{amssymb}
\usepackage{bm}
\usepackage[shortlabels]{enumitem}
\newtheorem{lemma}{Lemma}
\newtheorem{proposition}{Proposition}
\DeclareMathOperator{\cl}{cl}
\newcommand{\forkindep}[1][]{%
  \mathrel{
    \mathop{
      \vcenter{
        \hbox{\oalign{\noalign{\kern-.3ex}\hfil$\vert$\hfil\cr
              \noalign{\kern-.7ex}
              $\smile$\cr\noalign{\kern-.3ex}}}
      }
    }\displaylimits_{#1}
  }
}
\newcommand{\isom}{\cong}
\title{Classification strength of Polish groups: involving $S_\infty$}
\author{Shaun Allison}
\date{Source: arXiv:2304.00139v2 (CC BY 4.0)}
\begin{document}
\maketitle

\noindent\emph{Source and licence.} Classification strength of Polish groups: involving $S_\infty$. Version arXiv:2304.00139v2 (CC BY 4.0), distributed by arXiv under the Creative Commons Attribution 4.0 International licence, http://creativecommons.org/licenses/by/4.0/, as recorded in the arXiv licence metadata for this exact version and shown on the abstract page https://arxiv.org/abs/2304.00139v2. Full licence notice: \texttt{licenses/arxiv-2304.00139-CC-BY-4.0.txt}.\par\smallskip

\noindent\textit{Editorial scope.} Counted unit: the article's proposition prop:cl\_equivalence (source0.tex lines 788--801). The article's complete original proof of it is delivered with it (source0.tex lines 803--897, one \texttt{proof} environment of 95 lines). No mathematics is written, repaired or re-derived: the statement and the proof are the article's own lines, in the article's own notation, and no retained line is substituted or re-flowed.  Retained uncounted as the source-local context the proof needs: lines 767--769.  Nothing was omitted from any retained span, so the package carries no omission record.  No retained line contains a citation, so the wrapper carries no bibliography and no bibliography entry.  No retained line contains a figure environment, an image-inclusion command or drawing code, so no image is delivered and the package declares no asset dependency.  Editorial formatting only, in addition: an AMS \texttt{amsart} preamble that restates the article's own declarations and macros used by the retained lines, namely the environments \texttt{lemma}, \texttt{proposition} on separate counters, together with the standard packages those lines need.  The retained environments print in this document as Lemma 1, Proposition 1 in order of appearance, whereas the article prints the same items with the numbering of its own full document.  The complete native source of the article as distributed in the arXiv e-print for this exact version is bundled unchanged as \texttt{sources/139533/source0.tex}.
\begin{lemma}\label{lem:one_side}
Suppose $P \curvearrowright I$ and $\cl$ is an invariant closure operator on $I$. Then for any $a, b \in I$ and $C \subseteq I$ finite, if for every $a' \isom_{C} a$ we have $a' \in \cl(bC)$, then for every $a' \isom_{C} a$ and $b' \isom_{C} b$ we have $a' \in \cl(b' C)$.
\end{lemma}

\begin{proposition}\label{prop:cl_equivalence}
Suppose $\cl$ is an invariant finite-character closure operator on $P \curvearrowright I$ with derived independence relation $\forkindep$. The following are equivalent:
\begin{enumerate}
    \item $\cl$ is disjointifying;
    \item for any finite $C \subseteq I$ and $B \subseteq I$ with $C \subseteq B$ and $a \in I$, there is some $a' \isom_C a$ such that $a'C \forkindep[C] B$; and
    \item for any finite $C \subseteq I$ and $a, b \in I$, there is some $a' \isom_C a$ such that $a'C \forkindep[C] bC$;
    \item for any finite $C \subseteq I$ and $a, b \in I$ with $a \not\in \cl(C)$, the following both hold:
    \begin{enumerate}
        \item there is some $a' \isom_C a$ such that $a'C \forkindep[C] aC$; and
        \item there is some $a' \isom_C a$ such that $a' \not\in \cl(bC)$.
    \end{enumerate}
\end{enumerate}

\end{proposition}

\begin{proof}
Clearly (1) implies (2), (2) implies (3), and (3) implies (4).
However showing (4) implies (1) is not easy, and instead we show each step backwards. In order we show (3) implies (2), (2) implies (1), then (4) implies (3).

We prove (3) implies (2).
We will induct on $|B \setminus \cl(C)|$.
For the base case, let $B, C \subseteq I$ finite with $C \subseteq B$ and $B \subseteq \cl(C)$.
Then for any $a \in I$ by existence we have $aC \forkindep[C] B$.

For the induction step, let $B, C \subseteq I$ be finite and assume $|B \setminus \cl(C)| > 0$ and that the claim holds for lesser values of $|B \setminus \cl(C)|$.
Let $a \in I$ and fix some $b_0 \in B \setminus \cl(C)$ minimal in $B$ over $C$.
By assumption find $a' \isom_C a$ such that $a'C \forkindep[C] b_0C$. 
By the induction hypothesis, find some $a'' \isom_{b_0C} a'$ as witnessed by some $\pi \in P$ such that $a''C \forkindep[b_0C] B$.
Clearly $a'' \isom_C a$.
By invariance with respect to $\pi$ we have $a''C \forkindep[C] b_0C$, and then by weak transitivity we have $a''C \forkindep[C] B$ as desired.

The proof that (2) implies (1) is similar.
We induct on $|A \setminus \cl(C)|$. 
If this is 0, then as before the conclusion follows by existence.
Otherwise, assume $|A \setminus \cl(C)| > 0$ and the claim holds for lesser values.
Let $a_0 \in A \setminus \cl(C)$ minimal in $A$ over $C$.
Find $a_0' \isom_C a_0$ such that $a_0' \forkindep[C] B$, as witnessed by some $\pi \in P$.
Define $A' := \pi[A]$ so that $A' \isom_C A$.
By invariance of $\cl$, we have $a_0' \in A'$ is minimal in $A'$ over $C$.
By the induction hypothesis, find $A'' \isom_{a_0'C} A'$ such that $A'' \forkindep[a_0'C] B$.
Observe that $a_0'$ is minimal in $A''$ over $C$ as well, by invariance of $\cl$.
By weak transitivity we have $A'' \forkindep[C] B$ as desired.

Finally we see that (4) implies (3). 
We prove this in a few steps via the following intermediate statements
\begin{enumerate}
    \item[(3*)] Given $a, C$ suppose there exists some $b \not\in \cl(C)$ such that for every $b' \isom_C b$, either $a \in \cl(b'C)$ or $b' \in \cl(aC)$. Then $a \in \cl(C)$.
    \item[(4*)] Given $a, C$ suppose either of the following holds:
    \begin{enumerate}
        \item[(a)] there exists some $b$ such that for every $b' \isom_C b$, $a \in \cl(b'C)$; or
        \item[(b)] for any $a' \isom_C a$ either $a \in \cl(a'C)$ or $a' \in \cl(aC)$.
    \end{enumerate}
    Then $a \in \cl(C)$.
\end{enumerate}
We show that (4) implies (4*) and (4*) implies (3*) and then (3*) implies (3).

Let's show that (4) implies (4*).
Let $a \in I$ and $C \subseteq I$ finite.
First let's suppose (4*).(a) holds, i.e. we have some $b \in I$ such that for every $b' \isom_C b$ we have $a \in \cl(b'C)$.
If we assume for the sake of contradiction that $a \not\in \cl(C)$ then by (4).(b) there is some $a' \isom_C a$ as witnessed by some $\pi \in P$ such that $a' \not\in \cl(bC)$.
Letting $b' := \pi^{-1}(b)$ we have by invariance of $\cl$ that $a \not\in \cl(b'C)$ where $b' \isom_C b$, which contradicts our choice of $b$.
Next, let's suppose (4*).(b) holds, i.e. for any $a' \isom_C a$ either $a \in \cl(a'C)$ or $a' \in \cl(aC)$.
If we assume for the sake of contradiction that $a \not\in \cl(C)$ then by (4).(a) there is some $a' \isom_C a$ such that $a'C \forkindep[C] aC$, or in other words
\[ a'C \cap \cl(aC) \subseteq \cl(C) \quad \text{and} \quad aC \cap \cl(a'C) \subseteq \cl(C). \]
However if $a \in \cl(a'C)$ then by the right inclusion we have $a \in \cl(C)$ which is a contradiction, and if $a' \in \cl(aC)$ then by the left inclusion we have $a' \in \cl(C)$ and thus $a \in \cl(C)$ by invariance of $\cl$, which is also a contradiction.

Next we show that (4*) implies (3*)
Suppose we're given $a, C$ and some $b \not\in \cl(C)$ such that for every $b' \isom_C b$, either $a \in \cl(b'C)$ or $b' \in \cl(aC)$. 
Assume for contradiction that $a \not\in \cl(C)$. 

We in fact have that for every $a' \isom_C a$ and $b' \isom_C b$, either $a' \in \cl(b'C)$ or $b' \in \cl(a'C)$.
Indeed, let $\pi$ witness $a' \isom_C a$ and let $b'' := \pi(b')$.
Then we have $b'' \isom_C b$ and so either $a \in \cl(b''C)$ or $b'' \in \cl(aC)$.
In the first case by invariance of $\cl$ with respect to $\pi^{-1}$ we have $a' \in \cl(b'C)$ and in the second case for the same reason we have $b' \in \cl(a'C)$ as desired.

Define a relation $\le$ on the set of all $a'$ with $a' \isom_C a$ by
\[ a' \le a'' \quad \text{iff} \quad \forall b' \isom_C b, \; (a'' \in \cl(b'C) \rightarrow a' \in \cl(b'C)). \]
We argue that this is in fact a linear preorder.
Reflexivity and transitivity of $\le$ are clear. 
To verify trichotomy, suppose for contradiction there exist some $a_0, a_1 \isom_C a$ with $a_0 \not\le a_1$ and $a_1 \not\le a_0$ or in other words, there are some $b_0, b_1 \isom_C b$ with 
\[\text{[i]} \; a_1 \in \cl(b_0C) \quad \text{and} \quad \text{[ii]} \; a_0 \not\in \cl(b_0C)\]
and 
\[\text{[iii]} \; a_1 \not\in \cl(b_1C) \quad \text{and} \quad \text{[iv]} \; a_0 \in \cl(b_1C).\]
By [iii], we have $b_1 \in \cl(a_1C)$ and thus by [iv] we have $a_0 \in \cl(a_1C)$. 
But that implies $a_0 \in \cl(b_0C)$ by [i], which contradicts [ii].

Now we show that $a' \le a''$ implies that $a' \in \cl(a''C)$.
Since either $a \le a'$ or $a' \le a$ for every $a' \isom_C a$, we would be done because (4*).(b) would hold and we would have $a \in \cl(C)$, a contradiction.

Assume $a' \le a''$.
We may find some $b' \isom_C b$ such that $a'' \in \cl(b'C)$.
Indeed, otherwise we would have $b' \in \cl(a''C)$ for every $b' \isom_C b$.
Thus $b' \in \cl(a'''C)$ for every $a''' \isom_C a$ by Lemma \ref{lem:one_side}, allowing us to conclude $b \in \cl(C)$ by (4*).(a), a contradiction.

Then of course by invariance we have $a'' \in \cl(b''C)$ for every $b'' \isom_{a''C} b'$ and since $a' \le a''$ we thus have $a' \in \cl(b''C)$ for every $b'' \isom_{a''C} b'$.
By (4*).(a), we can conclude $a' \in \cl(a''C)$ as desired.

Finally, we show that (3*) implies (3).
Assume (3*) and let finite $C \subseteq I$ and $a \in I$ be arbitrary.
Assume for the sake of contradiction that there is some $b \in I$ such that for every $a' \isom_C a$, we have $a'C \not\forkindep[C] bC$.
We must have $a, b \not\in \cl(C)$ otherwise $aC \forkindep[C] bC$ would follow by existence and symmetry.
We claim that
\[ \forall b' \isom_C b, \; a \in \cl(b'C) \quad \text{or} \quad b' \in \cl(aC). \]
This suffices as by (3*) we would conclude $a \in \cl(C)$ which would be a contradiction.
To that end, fix arbitrary $b' \isom_C b$ as witnessed by some $\pi \in P$.
Define $a' := \pi(a)$ in which case we have $a' \isom_C a$ and so $a'C \not\forkindep[C] bC$ by assumption.
By invariance under $\pi$ we have $aC \not\forkindep[C] b'C$ and thus we have
\[ aC \cap \cl(b'C) \not\subseteq \cl(C) \quad \text{or} \quad b'C \cap \cl(aC) \not\subseteq \cl(C). \]
In the first case, we conclude that $a \in \cl(b'C)$ and in the second case we conclude $b' \in \cl(aC)$ as desired.
\end{proof}

\end{document}
